\documentclass[conference]{IEEEtran}
\IEEEoverridecommandlockouts
% The preceding line is only needed to identify funding in the first footnote.
% If that is unneeded, please comment it out.
% Template version as of 6/27/2024

\usepackage{cite}
\usepackage{amsmath,amssymb,amsfonts}
\usepackage{algorithmic}
\usepackage{graphicx}
\usepackage{textcomp}
\usepackage{xcolor}
\usepackage{booktabs}   % nicer tables (\toprule, \midrule, \bottomrule)
\usepackage{siunitx}    % consistent units, e.g. \SI{5}{\centi\meter}
\usepackage[hidelinks]{hyperref}

\graphicspath{{figures/}}

% Red TODO markers -- search for \todo before submission; none may remain.
\newcommand{\todo}[1]{\textcolor{red}{[TODO: #1]}}
% Checkmark / dash for the comparison table
\newcommand{\cmark}{$\checkmark$}
\newcommand{\xmark}{--}

\begin{document}

% CRITICAL: Do not use symbols, special characters, footnotes, or math in
% the paper title or abstract.
\title{SCOOPS: A Holistic Multi-Agent Dataset for Sensing, Communication, and Cooperative Perception\\
% Delete the \thanks line if there is no funding to acknowledge.
\thanks{Identify applicable funding agency here. If none, delete this.}
}

\author{%
\IEEEauthorblockN{Given Name Surname}
\IEEEauthorblockA{\textit{Department of Electrical and Computer Engineering} \\
\textit{National Yang Ming Chiao Tung University}\\
Hsinchu, Taiwan \\
jorich.ee14@nycu.edu.tw}
% Add co-authors by repeating the block below, separated by \and
%\and
%\IEEEauthorblockN{2\textsuperscript{nd} Given Name Surname}
%\IEEEauthorblockA{\textit{dept. name of organization (of Aff.)} \\
%\textit{name of organization (of Aff.)}\\
%City, Country \\
%email address or ORCID}
}

\maketitle

\begin{abstract}
Cooperative robotic systems must jointly solve perception, sensing, and
communication, yet existing datasets capture at most two of these layers at
once. We present SCOOPS, a holistic multi-agent, multi-view, multi-sensor
dataset that spans all three: synchronized camera, radar, LiDAR, and RTK-GNSS sensing;
cooperative perception annotations; and measured, replayable wireless channel
traces. Each platform is spatially calibrated and temporally aligned, enabling
tasks across single-agent and collaborative scenarios and modalities. We
provide ground-truth poses, TSDF reconstruction supervision, RSSI channel
traces, and object annotations, together with benchmarks for cooperative
perception under clean and actual channels, point cloud reconstruction, and
cross-modal localization. \todo{One sentence with the headline result, e.g.
``Cooperative perception improves mAP by X points over single-agent, and the
logged channel reduces this gain by Y.''} The dataset and development kit are
publicly available at \todo{URL/DOI}.
\end{abstract}

\begin{IEEEkeywords}
multi-agent systems, cooperative perception, collaborative SLAM, integrated
sensing and communication, dataset, benchmark
\end{IEEEkeywords}

%=======================================================================
\section{Introduction}
\label{sec:introduction}
% Goal (from outline): establish that cooperative robotics needs all three
% layers jointly, no dataset provides that, and SCOOPS does.
Cooperative robotic systems---fleets of vehicles, drone swarms, and
multi-robot teams---depend on three tightly coupled layers: \emph{sensing}
for localization and mapping, \emph{perception} for scene understanding, and
\emph{communication} for sharing information between agents. Progress on any
one layer in isolation does not transfer to deployed systems, because the
layers interact: perception quality depends on what the channel can deliver,
and localization accuracy bounds how observations from different agents can
be fused. \todo{Add one or two motivating application sentences.}

Existing datasets, however, cover these layers only partially.
Collaborative SLAM datasets provide rich synchronized sensing but treat
communication implicitly \cite{s3e,kimeramulti}; cooperative perception
benchmarks provide labels but idealize or simulate the channel
\cite{opv2v,dairv2x,v2v4real}; and joint sensing--communication datasets log
the real channel but offer no cooperative scene understanding
\cite{deepsense6g}. As a result, questions that span layers---how much does
a real, lossy channel erode cooperative perception gains? which modality
localizes best under communication constraints?---cannot be answered with
any single existing dataset.

We close this gap with \textbf{SCOOPS} (\textbf{S}ensing,
\textbf{C}ommunication, c\textbf{OO}perative \textbf{P}erception
\textbf{S}uite), a holistic multi-agent, multi-view, multi-sensor dataset
that jointly captures all three layers plus dense reconstruction
supervision. The main contributions of this paper are:
\begin{itemize}
    \item \textbf{A tri-layer dataset.} Synchronized camera, radar, LiDAR,
    and RTK-GNSS streams from multiple spatially calibrated and temporally
    aligned agents, with measured wireless channel traces recorded alongside
    the sensing data.
    \item \textbf{Comprehensive ground truth.} Per-environment pose ground
    truth, TSDF-based reconstruction supervision, replayable RSSI channel
    traces, and cooperative perception annotations.
    \item \textbf{Benchmarks and baselines.} Task definitions, splits, and
    baseline results for cooperative perception under clean and actual
    channels, point cloud reconstruction, and localization from camera,
    radar, and CSI.
    \todo{Adjust the contribution list once results are final.}
\end{itemize}

%=======================================================================
\section{Related Work}
\label{sec:related}
% Theme (from outline): each family of related work lacks one dimension --
% SLAM, perception, sensing, or communication.

\subsection{Collaborative SLAM Datasets}
% Sensing-centric: rich sensing, but targets localization/mapping and
% treats communication only implicitly.
Multi-robot SLAM datasets such as S3E \cite{s3e} and systems such as
Kimera-Multi \cite{kimeramulti} provide rich, synchronized multi-agent
sensing aimed at collaborative localization and mapping. However, they
target the sensing layer alone: inter-agent communication is either assumed
ideal or logged only as a system artifact, and cooperative perception is not
considered. \todo{Add 2--3 more datasets (e.g., GRACO, FordAV, EuRoC-style
multi-robot sets) and one sentence each.}

\subsection{Cooperative Perception and Communication Datasets}
% Perception-centric: labels exist, but communication is idealized or
% simulated, not measured.
Cooperative perception benchmarks provide multi-agent perception labels,
either in simulation (OPV2V \cite{opv2v}) or from real vehicles and
infrastructure (DAIR-V2X \cite{dairv2x}, V2V4Real \cite{v2v4real}).
Communication-efficient methods such as Where2comm \cite{where2comm} are
evaluated on these benchmarks under idealized bandwidth models. In all
cases the channel itself is simulated or abstracted away, so the effect of
a real, time-varying link on cooperation cannot be measured.
\todo{Verify and extend the paper list.}

\subsection{Joint Sensing--Communication Datasets}
% Communication-centric: real channel is logged, but no cooperative scene
% understanding.
Integrated sensing and communication datasets such as DeepSense~6G
\cite{deepsense6g} log real wireless channels co-registered with sensing
modalities. These datasets ground the communication layer in measurement,
but they are single-link and task-agnostic: they contain no multi-agent
cooperation and no scene-understanding labels. \todo{Add further ISAC/CSI
datasets and one sentence each.}

\subsection{3D Reconstruction Benchmarks}
% Single-agent, posed imagery only; no cooperative or TSDF-supervised
% configuration.
Reconstruction benchmarks such as ScanNet \cite{scannet} and Tanks and
Temples \cite{tanksandtemples} supervise dense geometry from posed
single-agent imagery. They provide no multi-agent configuration and no
communication dimension, so cooperative reconstruction---where partial
observations must be fused across a constrained link---cannot be studied.

\subsection{Positioning of SCOOPS}
Table~\ref{tab:comparison} summarizes the landscape. Each prior family
covers at most two of the three layers; SCOOPS is, to our knowledge, the
only dataset spanning synchronized multi-agent sensing, cooperative
perception labels, and measured communication, together with dense
reconstruction supervision.

\begin{table}[htbp]
\caption{Comparison of SCOOPS With Existing Datasets}
\begin{center}
\begin{tabular}{lcccc}
\toprule
\textbf{Dataset} & \textbf{Multi-} & \textbf{Percep.} & \textbf{Measured} & \textbf{Recon.} \\
 & \textbf{agent} & \textbf{labels} & \textbf{comm.} & \textbf{superv.} \\
\midrule
S3E \cite{s3e}                     & \cmark & \xmark & \xmark & \xmark \\
OPV2V \cite{opv2v}                 & \cmark & \cmark & \xmark & \xmark \\
DAIR-V2X \cite{dairv2x}            & \cmark & \cmark & \xmark & \xmark \\
DeepSense 6G \cite{deepsense6g}    & \xmark & \xmark & \cmark & \xmark \\
ScanNet \cite{scannet}             & \xmark & \cmark & \xmark & \cmark \\
\textbf{SCOOPS (ours)}             & \cmark & \cmark & \cmark & \cmark \\
\bottomrule
\end{tabular}
\label{tab:comparison}
\end{center}
\end{table}

%=======================================================================
\section{The SCOOPS Dataset}
\label{sec:dataset}
The SCOOPS system focuses on bridging the gap between sensing and
communication for SLAM and perception tasks. Each platform is spatially
calibrated and temporally aligned to provide multi-agent, multi-view,
multi-sensor information, enabling a wide range of tasks across single-agent
and collaborative scenarios and across modalities.
\todo{Add a system-overview figure (Fig.~X) showing agents, sensors, and
the communication link.}

\subsection{Sensing Suite}
% Specify the synchronized payload per agent.
Each agent carries a synchronized payload of camera, radar, LiDAR, network
interface card (NIC), and RTK-GNSS. Table~\ref{tab:sensors} lists the
sensor models, rates, and resolutions.
\todo{Fill in exact sensor models, rates, FoV, and mounting details; state
whether IR/IMU streams are included.}

\begin{table}[htbp]
\caption{Per-Agent Sensor Suite}
\begin{center}
\begin{tabular}{llll}
\toprule
\textbf{Sensor} & \textbf{Model} & \textbf{Rate} & \textbf{Resolution / Notes} \\
\midrule
Camera    & \todo{model} & \todo{Hz} & \todo{resolution, FoV} \\
Radar     & \todo{model} & \todo{Hz} & \todo{range, FoV} \\
LiDAR     & \todo{model} & \todo{Hz} & \todo{channels, range} \\
NIC       & \todo{model} & \todo{Hz} & \todo{band, protocol} \\
RTK-GNSS  & \todo{model} & \todo{Hz} & \todo{accuracy} \\
\bottomrule
\end{tabular}
\label{tab:sensors}
\end{center}
\end{table}

\subsection{Spatial Calibration}
\subsubsection{Agent}
\todo{Describe intrinsic and extrinsic calibration of each agent's sensors
(e.g., ChArUco-based camera calibration, radar--camera and LiDAR--camera
extrinsics) and report residual errors.}

\subsubsection{Infrastructure}
\todo{Describe how infrastructure-mounted sensors are calibrated into the
common frame and report accuracy.}

\subsection{Temporal Calibration}
\subsubsection{NTP Synchronization}
All agents are synchronized over NTP. \todo{Report achieved clock offset
statistics (mean/max), and describe hardware triggering if any.}

\subsubsection{Motion Compensation and Fusion Error Bounds}
% Report worst-case intra-frame spatial error at max agent speed.
\todo{Derive/report the worst-case intra-frame spatial error at the maximum
agent speed, combining clock offset and sensor exposure/sweep time, e.g.
``at \SI{2}{\meter\per\second} and a worst-case offset of X ms, the bound
is Y cm.''}

%=======================================================================
\section{Dataset Design and Collection}
\label{sec:collection}

\subsection{Environments}
% Indoor/outdoor sites chosen for viewpoint, appearance, and radio diversity.
\todo{Describe the indoor and outdoor sites and why they were chosen:
viewpoint diversity, appearance diversity, and radio-propagation diversity
(LoS/NLoS, multipath). Add a site-overview figure.}

\subsection{Collection Protocol}
\todo{Describe trajectories, agent roles, number of runs per site, operator
procedure, and any scripted interactions between agents.}

\subsection{Conditions Captured}
% Inventory the deliberate difficulty and totals.
\todo{Inventory the deliberately captured difficulty: low-light, sensor
dropout, packet loss, occlusion. Report totals: hours, distance, frames,
annotated instances.}

%=======================================================================
\section{Ground Truth and Annotations}
\label{sec:groundtruth}

\subsection{Pose Ground Truth}
\todo{Report per-environment pose ground truth source (RTK-GNSS outdoors,
motion capture / map-based indoors?) and its accuracy.}

\subsection{Reconstruction Supervision}
Dense reconstruction ground truth is produced with GLIM \cite{glim} in
ROS~2, providing globally consistent maps used as TSDF supervision.
\todo{Describe the mapping pipeline settings, map resolution, and how map
quality was validated.}

\subsection{Communication Ground Truth}
% Channel traces as replayable/resamplable ground truth.
Wireless channel measurements are logged as time-stamped RSSI traces
aligned with the sensing streams, so the channel can be replayed or
resampled as ground truth for communication-aware experiments.
\todo{State the measurement setup (band, packet rate), what is logged
(RSSI, CSI?, throughput, loss), and how traces align with sensor frames.}

\subsection{Perception Annotations and Statistics}
\todo{Detail label types (3D boxes? semantic masks?), the annotation
protocol and QA process, and class statistics (instances per class, per
environment). Add a statistics figure.}

%=======================================================================
\section{Data Format, Structure, and Access}
\label{sec:format}
% rosbag, recorded + synced timestamps, devkit, license, DOI, anonymization.
Recordings are distributed as ROS~2 bags preserving both the originally
recorded timestamps and the synchronized timestamps, alongside a
development kit for parsing, calibration lookup, and benchmark evaluation.
\todo{Specify the directory layout, per-topic message types, devkit
features, license, DOI, and the anonymization procedure (faces/plates).}

%=======================================================================
\section{Benchmark Tasks and Baselines}
\label{sec:benchmarks}

\subsection{Task Definitions}
We define four benchmark tasks: (i)~cooperative point cloud
reconstruction, (ii)~cooperative object detection, (iii)~cross-view
localization from camera, radar, and CSI, and (iv)~communication-efficient
perception under the logged channel.
\todo{Give each task a precise input/output definition.}

\subsection{Splits and Metrics}
\todo{Define the within-environment and leave-one-environment-out
protocols, and the metric for each task (e.g., mAP for detection, Chamfer
distance / F-score for reconstruction, ATE for localization, and
accuracy-vs-bandwidth curves for comm-efficient perception).}

\subsection{Baselines}
\todo{List baselines per task: cooperative perception (clean channel vs.
actual logged channel), point cloud reconstruction, and localization from
camera, radar, and CSI. Cite each baseline method.}

\subsection{Main Results}
% Headline cooperative-vs-single-agent and bandwidth trade-off numbers.
\todo{Report the headline numbers: cooperative vs. single-agent
performance, and the bandwidth/accuracy trade-off under the real channel.
Add the main results table(s).}

\subsection{Ablations}
% Isolate each layer, including idealized vs. logged channel.
\todo{Isolate the contribution of each layer: sensing modalities on/off,
idealized vs. logged channel, number of cooperating agents.}

\subsection{Qualitative Results}
\todo{Show cooperative vs. single-agent reconstruction side by side, and a
failure case under channel degradation.}

%=======================================================================
\section{Discussion}
\label{sec:discussion}

\subsection{Findings}
\todo{State where cooperation helps, where it does not, and how conclusions
shift when the idealized channel is replaced by the logged one.}

\subsection{Limitations}
\todo{Bound the dataset plainly: scale, number of agents, ground-truth
variance, and radio specificity (one band/protocol).}

\subsection{Broader Impact and Ethics}
\todo{Address dual-use considerations, licensing terms, and the
anonymization of captured people and vehicles.}

%=======================================================================
\section{Conclusion and Future Work}
\label{sec:conclusion}
We presented SCOOPS, a holistic multi-agent dataset that jointly captures
synchronized sensing, cooperative perception labels, and measured
communication, with dense reconstruction supervision and benchmarks across
the three layers. \todo{One sentence of headline findings.} Future work
includes scaling to larger teams, adding modalities, and richer channel
models.

\section*{Acknowledgment}
% Optional; delete if unused. Sponsor acknowledgments go in the unnumbered
% footnote on the first page (\thanks in \title).
\todo{Add acknowledgments or delete this section.}

\bibliographystyle{IEEEtran}
\bibliography{references}

\end{document}
